% ==================================================================
%  CorrectnessDetails.tex  --  Appendix A
%  Details of Section 3.2.4 (Section3/Correctness.tex): derivation of the
%  precisions, proof of Theorem 1, sign bootstrapping and cleaning, and
%  the error of the masks and of Max.
%  Loaded from main.tex after the references, under \appendices.
% ==================================================================

\section{Details of the Correctness Analysis}
\label{app:correct}

For a precision $\alpha$, the composite approximation of
Section~\ref{sec:comparison} returns a value within $2^{1-\alpha}$ of the
sign of its input $x$ whenever $|x| \ge 2^{-\alpha}$; an input inside the
dead zone $(-2^{-\alpha}, 2^{-\alpha})$ may give any value in $[-1,1]$.
A comparison is correct when the inputs that decide the output lie
outside the dead zone, and it is accurate when the value it returns there
is close to $\pm 1$. This appendix derives the precisions of
Section~\ref{sec:correct}, proves Theorem~\ref{thm:sandwich}, and
describes how the server brings the outputs closer to $\pm 1$, and the
masks closer to $0$ and $1$.

\subsection{Dead Zone in Neighbor Check}
\label{app:neighbor}

Nothing in the public parameters bounds how close a squared distance can
come to $\varepsilon_n^2$, so the server uses a precision $\alpha_d$
large enough to make the dead zone narrow and chooses on which side of
$\varepsilon_n$ it lies. With the threshold at $\varepsilon_n^2$ itself,
the dead zone would lie on both sides of it, and a pair inside it could
be decided wrongly in either direction. A true neighbor whose input lies
just below $0$ may come out with an ambiguous value and be missed, a false
negative. A non-neighbor whose input lies just above $0$ may come out
with an ambiguous value and be partly counted, a false positive. Removing
false negatives matters more: a lost neighbor can turn a core point into a
non-core point or cut a chain of core points, and so split a cluster that
plaintext DBSCAN keeps whole. The server therefore shifts the threshold
by the half-width of the dead zone, measured in squared distance, to
$\theta$ of Eq.~\eqref{eq:theta}. The dead zone then covers only the
squared distances between $\varepsilon_n^2$ and $\varepsilon_n^2 + 2\mu$,
all of them on the non-neighbor side. No false negative remains, and a
false positive adds a value between $0$ and $1$ to a count. Since only
false positives occur, the count of every point lies between the
plaintext counts at the radii $\varepsilon_n$ and
$\varepsilon_{\mathrm{hi}}$, which is Eq.~\eqref{eq:sandwich}.

\subsection{Dead Zone in Core Check and Proof of Theorem~\ref{thm:sandwich}}
\label{app:core}

Core Check compares $\hat{c}_i$ with $\mathit{MinPts} - \tfrac{1}{2}$
(Section~\ref{sec:core}). When the count is an integer, the shifted count
is a half-integer, so
$|\hat{c}_i - (\mathit{MinPts} - \tfrac{1}{2})| \ge \tfrac{1}{2}$, and
the input of Eq.~\eqref{eq:coreinput} satisfies $|y_i| \ge 1/(2n)$. The
smallest precision whose dead zone excludes every such input is
$\alpha_c = \lceil \log_2 (2n) \rceil$, since then
$2^{-\alpha_c} \le 1/(2n)$. A count is not an integer only when a false
positive of Neighbor Check contributes to it. Eq.~\eqref{eq:sandwich}
bounds the effect of such pairs, and the core test turns that bound into
the subset relation of Theorem~\ref{thm:sandwich}.

\begin{IEEEproof}[Proof of Theorem~\ref{thm:sandwich}]
We ignore the error of $\mathsf{Sgn}$ outside its dead zone, which
Appendices~\ref{app:signboot} and~\ref{app:error} bound.
If $i \in \mathrm{Core}_{\varepsilon_n}$, Eq.~\eqref{eq:sandwich} gives
$\hat{c}_i \ge \mathrm{Nei}_{\varepsilon_n}(i) \ge \mathit{MinPts}$, so
$y_i \ge 1/(2n) \ge 2^{-\alpha_c}$ and $\mathsf{Core\_mask}[i] = 1$. If
$i \notin \mathrm{Core}_{\varepsilon_{\mathrm{hi}}}$, it gives
$\hat{c}_i \le \mathrm{Nei}_{\varepsilon_{\mathrm{hi}}}(i) \le
\mathit{MinPts} - 1$, so $y_i \le -1/(2n)$ and
$\mathsf{Core\_mask}[i] = 0$. If no distance lies strictly between
$\varepsilon_n$ and $\varepsilon_{\mathrm{hi}}$, then
$\mathrm{Nei}_{\varepsilon_n}(i) = \mathrm{Nei}_{\varepsilon_{\mathrm{hi}}}(i)$
for every $i$, and the two bounds meet.
\end{IEEEproof}

\subsection{Dead Zone in Label Propagation}
\label{app:lp}

Every initial label is a multiple of $1/n$, and so is the $0$ of a masked
slot. Apart from its error, $\mathsf{Max}$ returns one of its two inputs,
so every label that $\mathsf{Max}$ compares is an initial label or $0$,
and two labels that differ do so by at least $1/n$. After the division by
$1.1$ in Eq.~\eqref{eq:step}, the input of $\mathsf{Sgn}$ then has
magnitude at least $1/(1.1\,n)$, and the smallest precision that keeps it
outside the dead zone is $\alpha_\ell = \lceil \log_2 (1.1\,n) \rceil$,
since then $2^{-\alpha_\ell} \le 1/(1.1\,n)$. Two equal labels give the
input $0$, inside the dead zone, but $\mathsf{Max}$ then returns their
common value whatever the sign (Section~\ref{sec:lp}). No input of the
maximum can therefore lead it to a wrong label.

\subsection{Sign Bootstrapping and Cleaning}
\label{app:signboot}

Outside the dead zone the composite polynomial returns a value within
$2^{1-\alpha}$ of $\pm 1$. The server pushes its outputs further toward
$\pm 1$, and with them the masks further toward $0$ and $1$, in two
ways.

First, following Hong et al.~\cite{hong2026white}, it refreshes the
ciphertext with sign bootstrapping after every component of the
composite polynomial, the last one included. Sign bootstrapping evaluates
a sine function in place of the modular reduction of ordinary
bootstrapping. The sine removes the multiple of the modulus that
bootstrapping introduces, and its extrema are placed at the two signs, so
a value near $\pm 1$ is drawn toward $\pm 1$: an input within $\tau$ of
$z \in \{-1, 1\}$ comes out within $(\pi^2/8)\,\tau^2$ of
$z$~\cite[Thm.~1]{hong2026white}. The bootstrapping that the composite
polynomial needs between its components thus also shrinks the error
quadratically, without an extra level.

Second, after the last sign bootstrapping the server applies one cleaning
polynomial. The signs that enter $\mathsf{Max}$ are cleaned with
$g(x) = \tfrac{3}{2}x - \tfrac{1}{2}x^3$, which
fixes $\pm 1$, and the masks $\mathsf{adj}_k$ and $\mathsf{Core\_mask}$
with $h(x) = 3x^2 - 2x^3$, which fixes $0$ and
$1$. Both polynomials have zero derivative at the points they fix, so an
error $\eta$ becomes $O(\eta^2)$; for instance,
$h(1-\eta) = 1 - 3\eta^2 + 2\eta^3$ and $h(\eta) = 3\eta^2 - 2\eta^3$.

\subsection{Halving the Error}
\label{app:error}

The masks and the maximum use $\mathsf{Sgn}$ only through
$\tfrac{1}{2}(1 \pm \mathsf{Sgn})$, in Eqs.~\eqref{eq:adj},
\eqref{eq:core} and~\eqref{eq:step}. If $\mathsf{Sgn}$ is within $\eta$
of $\pm 1$, then $\tfrac{1}{2}(1 \pm \mathsf{Sgn})$ is within $\eta/2$ of
$0$ or $1$. The mapping thus halves the error, from $2^{1-\alpha}$ to
$2^{-\alpha}$ for the bare approximation, which is the precision $\alpha$
of a comparison in~\cite{lee2021minimax}. In the maximum of
Eq.~\eqref{eq:max} the error of $\mathsf{ste}$ is multiplied by
$|\mathbf{b} - \mathbf{a}| \le 1$, so one maximum deviates from the
larger label by at most $|\mathbf{b} - \mathbf{a}| \cdot
2^{-\alpha_\ell}$. These small deviations are why the client receives
real values rather than exact labels (Section~\ref{sec:post}).
